\uuid{Lg7G}
\titre{Racine}
\auteur{bodin ; hanani ; mzari}
\organisation{exo7}
\licence{CC BY-NC-SA 4.0}
\datecreate{2020-01-20}
\niveau{L1}
\module{Algèbre}
\chapitre{Polynôme, fraction rationnelle}
\sousChapitre{Racine, décomposition en facteurs irréductibles}
\difficulte{1}
\theme{Racines, factorisation, racine}
\exo7qcmid{lille-1.polynomes.019}

\contenu{
\begin{question}
\qclassification{105.03}
\qtags{motcle=racine}
Soit $P(X) = X^7 - 5 X^5 - 5 X^4 + 4 X^3 + 13 X^2 + 12 X + 4$.
\begin{answers}

    \good{$-1$ est une racine de $P$.}

    \bad{$0$ est une racine de $P$.}

    \bad{$1$ est une racine de $P$.}

    \good{$2$ est une racine de $P$.}
\end{answers}
\begin{explanations}
Calculer $P(\alpha)$. En fait $P(X) = (X-2)^2(X+1)^3(X^2+X+1)$.
\end{explanations}
\end{question}
}
